% Part: history
% Chapter: biographies 
% Section: emmy-noether

\documentclass[../../../include/open-logic-section]{subfiles}

\begin{document}

\olfileid[hi]{his}{bio}{noe}

\olsection{एमी नोएथर}

\olphoto{noether-emmy}{एमी नोएथर}

एमी नोएथर (\textsc{नर}-टर) का जन्म 23 मार्च 1882 को जर्मनी के एरलांगेन
में एक उच्च-मध्यवर्गीय विद्वान परिवार में हुआ। ``आधुनिक बीजगणित की जननी''
कही जाने वाली नोएथर ने महिलाओं की शिक्षा के मार्ग में भारी बाधाओं के बावजूद
गणित और भौतिकी दोनों में युगांतरकारी योगदान दिए। उस समय जर्मनी में लड़कियों
को कलाओं की शिक्षा दी जानी थी और उन्हें महाविद्यालय-तैयारी विद्यालयों में
जाने की अनुमति नहीं थी। फिर भी गॉटिंगेन और एरलांगेन विश्वविद्यालयों (जहाँ
उनके पिता गणित के प्राध्यापक थे) में कक्षाएँ सुनने के बाद नोएथर अंततः 1904
में एरलांगेन में विद्यार्थी के रूप में नामांकन कर सकीं, जब उसकी नीति बदलकर
महिला विद्यार्थियों को अनुमति दी गई। उन्होंने 1907 में गणित में डॉक्टरेट
प्राप्त की।

अपनी योग्यताओं के बावजूद नोएथर को कार्य-जीवन में बहुत विरोध झेलना पड़ा।
1908--1915 तक उन्होंने एरलांगेन में बिना वेतन पढ़ाया। इस दौरान उन्होंने उस
समय के विश्व के अग्रणी गणितज्ञों में से एक डेविड हिल्बर्ट का ध्यान खींचा;
हिल्बर्ट ने उन्हें गॉटिंगेन बुलाया। किंतु महिलाओं को प्राध्यापक पद पाने की
मनाही थी और वे केवल हिल्बर्ट के नाम से, फिर बिना वेतन, व्याख्यान दे सकीं।
इसी समय उन्होंने वह सिद्ध किया जिसे अब नोएथर का प्रमेय कहते हैं और जिसका
उपयोग आज भी सैद्धांतिक भौतिकी में होता है। अंततः 1919 में नोएथर को पढ़ाने का
अधिकार मिला। विश्वविद्यालय के सहकर्मियों के लगातार विरोध पर हिल्बर्ट की
प्रतिक्रिया कथित रूप से थी: ``सज्जनो, संकाय की सीनेट कोई स्नानागार नहीं
है।''

1920 के दशक के उत्तरार्ध में उन्होंने अमूर्त बीजगणित के कार्य पर ध्यान
केंद्रित किया और उनके योगदान ने इस क्षेत्र में क्रांति ला दी। अपने प्रमाणों
में वे प्रायः तथाकथित आरोही शृंखला प्रतिबंध का उपयोग करती थीं, जो कहता है कि
कुछ समुच्चयों की कोई अनंत कड़ाई से बढ़ती शृंखला नहीं होती। उदाहरणतः अब
नोएथरी वलय कहलाने वाली कुछ बीजीय संरचनाओं में गुण होता है कि आदर्शों का कोई
अनंत अनुक्रम $I_1 \subsetneq I_2 \subsetneq \dots$ नहीं होता। इस प्रतिबंध
को किसी भी आंशिक क्रम पर व्यापक बनाया जा सकता है (बीजगणित में इसका संबंध
उपसमुच्चय संबंध से क्रमित आदर्शों की विशेष स्थिति से है), और हम द्वैत अवरोही
शृंखला प्रतिबंध पर भी विचार कर सकते हैं, जिसमें किसी आंशिक क्रम का प्रत्येक
कड़ाई से \emph{घटता} अनुक्रम अंततः समाप्त होता है। यदि कोई आंशिक क्रम अवरोही
शृंखला प्रतिबंध को संतुष्ट करता है, तो इस क्रम के साथ आगमन का उसी तरह उपयोग
संभव है जैसे हम~$\Nat$ पर $<$ क्रम के साथ आगमन का उपयोग कर सकते हैं। ऐसे
क्रमों को \emph{सुस्थापित} या \emph{नोएथरी}, और संबद्ध प्रमाण सिद्धांत को
\emph{नोएथरी आगमन} कहते हैं।

नोएथर यहूदी थीं और 1933 में नाज़ियों के सत्ता में आने पर उन्हें पद से हटा
दिया गया। सौभाग्य से नोएथर पेनसिल्वेनिया के ब्रिन मॉर में अस्थायी पद के लिए
संयुक्त राज्य अमेरिका जा सकीं। वहाँ रहते हुए उन्होंने प्रिंसटन में भी
व्याख्यान दिए, यद्यपि उन्हें विश्वविद्यालय महिलाओं के प्रति स्वागतशील नहीं
लगा \citep[81]{Dick1981}। 1935 में गर्भाशय की गाँठ निकालने के लिए नोएथर का
ऑपरेशन हुआ। शल्यक्रिया से हुए संक्रमण के कारण उनकी मृत्यु हुई और उन्हें
ब्रिन मॉर में दफ़नाया गया।

\begin{reading} 
नोएथर की जीवनी के लिए \citet{Dick1981} देखें। पेरिमीटर इंस्टीट्यूट फॉर
थ्योरेटिकल फिज़िक्स के नोएथर के जीवन और प्रभाव पर व्याख्यान ऑनलाइन उपलब्ध
हैं \citep{Perimeter2015}। यदि पढ़ते-पढ़ते थक गए हों, तो \emph{Stuff You
Missed in History Class} में नोएथर के जीवन और प्रभाव पर पॉडकास्ट है
\citep{Frey2015}। नोएथर की संकलित रचनाएँ मूल जर्मन में उपलब्ध हैं
\citep{Noether1983}।
\end{reading}

\end{document}
